\documentclass[../main.tex]{subfiles}

\begin{document}

\chapter{V.5 附录自检清单}

\begin{motivation}
	\textbf{「行百里者半九十」}——完成只是开始，精进永无止境。\\
	在快速迭代的学术与工程实践中，我们常犯的错误是：只关注产出，忽视验证。
	
	\textbf{核心动机}：提供一个可重复、可执行的自检框架，帮助你在交付前系统性审视每一个关键环节，从结构完整性到假设一致性，确保你的工作经得起推敲，远离「完成了但不对」的陷阱。
\end{motivation}

\begin{logicmap}
	\textbf{本章逻辑结构：} \\
	\indent motivation (为什么需要自检) \\
	\quad $\Rightarrow$ 工程自检清单 (做什么) \\
	\quad \indent $\Rightarrow$ 数学模块映射检查 \\
	\quad \indent $\Rightarrow$ 假设声明检查 \\
	\quad \indent $\Rightarrow$ 分支展开检查 \\
	\quad \indent $\Rightarrow$ 正文对应检查 \\
	\quad \Rightarrow$ 改进迭代循环 (如何精进)
\end{logicmap}

\section{工程自检}

\textbf{读者疑问：}你有没有想过，为什么一个看似完善的理论框架，在具体实施时总会出现意想不到的问题？自检清单的价值不仅在于发现错误，更在于构建一种系统性思维——让每一处细节都经得起质询。

% 工程自检
\begin{itemize}
	\item 是否所有数学模块均有正文映射？\quad 是
	\item 是否引入未声明假设？\quad 否
	\item 是否存在未展开但已声明分支？\quad 是
	\item 是否保持与正文论断一一对应？\quad 是
\end{itemize}

\subsection{逻辑衔接}

本章的自检框架并非孤立的检查项堆砌，而是与全书的核心逻辑紧密相连：
\begin{itemize}
	\item \textbf{结构完整性}源自「结构产生效率」的原理——即使得每一模块都有其存在的必要位置；
	\item \textbf{假设显式化}呼应「中性选择」思想——不预设最优，只呈现所有条件；
	\item \textbf{分支展开}体现「人存原理」的视角——为何某些路径被保留而其他被排除；
	\item \textbf{正文对应}回归「语言的信息压缩」本质——每个论断都应能被追溯，避免信息损失。
\end{itemize}
自检不仅仅是在「找 bug」，更是在「验证逻辑一致性和理论自洽」——这是从「完成了」迈向「完成了且正确」的关键一步。

\end{document}
